\documentclass[12pt]{article}
\usepackage[T1]{fontenc}
\usepackage[margin=2.5cm]{geometry}
\usepackage{amsmath,amssymb,amsthm}
\usepackage{booktabs,array,longtable}
\usepackage{xcolor}
\usepackage{hyperref}
\hypersetup{colorlinks=true,linkcolor=blue,citecolor=blue,urlcolor=blue}
\linespread{1.00}
\newtheorem{theorem}{Theorem}[section]
\newtheorem{proposition}[theorem]{Proposition}
\newtheorem{lemma}[theorem]{Lemma}
\theoremstyle{definition}
\newtheorem{definition}[theorem]{Definition}
\theoremstyle{remark}
\newtheorem*{remark}{Remark}
\newcommand{\E}{\mathbb E}
\newcommand{\Var}{\operatorname{Var}}
\newcommand{\argmax}{\operatorname*{arg\,max}}
\newcommand{\todo}[1]{\par\smallskip\noindent\textcolor{red!65!black}{\textbf{Author decision:} #1}\par\smallskip}
\newcommand{\code}[1]{\nolinkurl{#1}}
\newcommand{\source}[1]{\par\smallskip\noindent{\small\raggedright\textit{Project evidence:} #1\par}\smallskip}
\title{A maximum likelihood estimate or a correct stop:\\certification for univariate distribution families}
\author{Author names pending}
\date{First draft, 8 November 2026}
\begin{document}
\maketitle
\todo{Confirm authorship, affiliations, and acknowledgements.}
\begin{abstract}
A converged optimizer need not return a maximum likelihood estimate. Even for
familiar univariate distributions, the likelihood can have several local
maxima, a finite unattained supremum, or an unbounded concentration path.
Numerical tolerances can turn a small finite improvement into a false report
of nonexistence. We describe a software contract that treats these outcomes
as part of the statistical calculation: return a globally certified maximizer
to a stated coefficient accuracy, or stop with a reason and, where available,
competing witnesses and likelihood bounds. Certificates are valid in exact
arithmetic conditional on computed comparisons, with documented practical
floating-point margins; they are not outward-rounded interval proofs.
We develop existence and uniqueness arguments for negative binomial,
Poisson--inverse Gaussian, Fr\'echet, Kumaraswamy and conditional Burr
likelihoods, and give explicit constructions separating multimodality,
nonidentification and equal global maxima. Zero-inflated counts connect the
computational problem to the unresolved remainder of Sampford's truncated
negative-binomial conjecture. Boundary likelihoods and local experiments show
why successful optimization and valid inference require separate guarantees.
The accompanying library covers 99 families, with certification coverage
reported as a dated, changing inventory rather than a universal accuracy claim.
\end{abstract}

\section{What should a likelihood fit return?}\label{sec:intro}
Maximum likelihood is a definition of an estimator, not an optimization
algorithm. For observations $x=(x_1,\ldots,x_n)$ and a declared parameter
space $\Theta$, write
\[
 \ell_x(\theta)=\sum_{i=1}^n\log f_\theta(x_i),\qquad
 L_x^*=\sup_{\theta\in\Theta}\ell_x(\theta).
\]
A finite vector with a small score need not attain $L_x^*$. A local maximum
can be inferior to another finite maximum or to a limiting model. A sequence
of improving estimates can leave every compact parameter set. A failed solve
can nevertheless return a plausible vector. None of these difficulties
requires regression, latent random effects or a high-dimensional model.

Our setting is independent observations from a univariate parametric family,
usually with one to five coefficients. The analyst may hold any supported
subset of coefficients. This qualification matters: holding a mean is not
holding a mode, and holding a standard deviation is not holding the scale of
an unstandardized kernel. A model that has no finite maximum with all
coefficients free may have a unique maximum after one coefficient is held.
Conversely, a conditional one-dimensional likelihood may already be bimodal.

Examples below make the distinction concrete. One paralogistic observation
at held scale gives three stationary points, with the larger likelihood far
from the first local maximum. A zero-inflated negative-binomial sample has a
proved finite optimum whose entire improvement over the zero-inflated Poisson
face is smaller than a historical face-comparison tolerance. A generalized
Pareto sample has two exactly equal global maxima, one uniform and the other
interior. A beta-negative-binomial likelihood has distinct coefficient
maximizers representing the same probability law. These are different
statistical outcomes and should be reported differently.

The organizing contract is therefore: \emph{return a certified global
maximizer to a stated accuracy, or a correct stop}. A stop can identify a
limiting model, prove unboundedness or nonattainment, retain unresolved
competitors, record an exhausted work budget, or report numerical failure.
An unresolved comparison does not justify selecting its preferred witness.
A budget stop does not establish nonexistence. The desired contract is
stronger than an optimizer's convergence flag and stronger than a local score
gate. Section~\ref{sec:standard} distinguishes this target from the coverage
of the present implementation.

The library is evidence for the contract rather than the definition of its
truth. Its header-only C++17 core fits without R, and its R and JavaScript
bindings transport outcomes. Family dossiers hold likelihood arguments,
parameter conventions, limiting faces and certificate scopes. Symbolic
densities, independent high-precision calculations and regression fixtures
supply different kinds of validation. A frozen regression value checks
stability; it cannot prove that the value is an MLE.

\subsection{Scope and prior work}
The ordinary negative-binomial existence problem has a long history
\cite{anscombe,levin,aragon,wang}. Multiple Cauchy likelihood roots are
classical \cite{reeds}; scalar generalized Pareto reductions are established
\cite{grimshaw}. Boundary asymptotics and nonregular endpoint models likewise
have established theories \cite{self,smith}. We use these connections to
locate mathematical dependencies, not to claim a new general theory of global
optimization.

All priority language in this draft follows the project's literature ledger.
A 40-query search completed on 8 November 2026 covered the 13 showcase rows
then requiring further search, and beta-prime limiting behavior. It included
Fr\'echet held-shape existence, skew exponential-power ties, beta-prime profile
ties, PIG uniqueness, truncated and inflated count likelihoods, half-$t$
profiles and constrained endpoint likelihoods. Earlier targeted searches
examined negative-binomial proofs, truncated-NB conjectures and inference.
A source was classified by what was actually obtained: theorem text,
preprint, abstract, metadata or failed retrieval. ``Not located after search''
means just that. It is not a novelty assertion or an exhaustive absence claim.

For example, Jones \cite{jones} gives the Kumaraswamy shape profile; the
strict-concavity conclusion below was not located in the inspected material.
Wingo's abstracts \cite{wingo83,wingo93} explicitly concern Burr modality and
uniqueness, leaving exact theorem overlap open. Willmot \cite{willmot} was
metadata-only in the PIG search; Carlson's relevant article could not be
retrieved. The PIG theorem therefore has a local mathematical proof but no
settled priority assessment. The Fr\'echet search read a fixed-location theorem,
a different problem from the held-shape theorem below. These limits belong
beside the claims they qualify.
\todo{Choose the final Statistical Science emphasis: a theorem-led account with software evidence, or a broader review with these results as examples. The current draft follows the former.}

\section{The certificate standard}\label{sec:standard}
\subsection{Statistical and numerical claims}
\begin{definition}[Conditional certificate]
For a declared family, coefficient holds, parameter domain and sample, a
conditional certificate consists of mathematical inequalities that imply the
reported global outcome, together with the computed comparisons on which
those inequalities depend. Numerical comparisons are accepted only with the
family's stated practical floating-point margins. The implication is an
exact-arithmetic statement conditional on those comparisons.
\end{definition}
This standard leaves an explicit arithmetic assumption. Elementary functions,
recurrence sums and restored coefficients are evaluated in floating-point
arithmetic. Practical allowances account for cancellation, accumulation and
coordinate transformations; uncertain signs cause a stop. The allowances are
not universal error theorems for the platform's mathematical library.
Outward-rounded interval arithmetic would remove some of this conditionality,
but is not claimed here. A 70-digit repetition corroborates a comparison; it
does not turn it into a theorem about all arithmetic errors.
For example, the held-scale Burr certificate uses $64\epsilon(1+\max_j|a_j|)$
for elementary operations and transcendental values, and
$64\epsilon(q+1)(1+\sum_j|a_j|)$ for a $q$-term sum, with binary64
machine epsilon $\epsilon=2^{-52}$. Other families use different documented
allowances adapted to their scores and conditioning. These practical
margins are assumptions of the conditional implication, not a uniform
outward-rounding guarantee.

A successful fit needs both a global argument and a coefficient enclosure.
Suppose strict concavity gives a unique root $\theta_*$, and signed comparisons
place it in $[a,b]$. Then a sufficiently narrow image of $[a,b]$ under every
public coefficient map supplies coefficient accuracy. A short interval in
$\log r$ can be useful; a short interval in a badly conditioned transformed
coordinate may say little about the requested coefficient. For a scalar
positive coefficient, a log bracket of width $w$ bounds the relative error of
its geometric midpoint by $e^{w/2}-1$.

With several maxima, a lower bound at a feasible point and upper bounds over
all competitors can certify a global winner. A Taylor envelope over a compact
cover is useful only after all exterior paths have been bounded. A local
Hessian can certify an isolated maximizer inside a box; it does not exclude
maxima outside the box. The limiting faces are part of that exclusion
problem. Restricting a grid to a convenient finite window supplies no such
argument.

The intended uniform accuracy control is relative \code{tol}, with target
default $10^{-10}$ in every free coefficient, and \code{iterlim} for work.
The source snapshot contains certificates with older, explicitly wider
coefficient bounds. For example, the held-scale Burr report certifies relative
coefficient-image widths $10^{-6}$; the PIG dossier describes a log-shape
bracket target at least $2\times10^{-7}$ before polishing. Polishing can
improve a representative without narrowing the global enclosure. Thus
``globally certified'' and ``certified to $10^{-10}$ in every coefficient''
are separate claims in this draft.
\todo{Reconcile the uniform $10^{-10}$ accuracy target with each delivered certificate before presenting that target as an achieved library-wide guarantee. Decide the treatment of coefficients equal or close to zero.}

\subsection{Typed outcomes and finite work}
A finite fit returns a verified representative and its certificate scope.
An attained domain endpoint, such as an NB point mass at zero, is a model
member. An infinite parameter corresponding to a Poisson law is a limit face,
not a finite NB estimate. A restricted estimate is labeled as constrained;
the restriction is not silently interpreted as a theorem about the
unrestricted family. Invalid observations are distinguished from valid data
whose likelihood has no maximum.

A tie stop means that separated competitors cannot be ordered at the stated
arithmetic resolution. It need not mean an exact mathematical tie. Exact ties
below are established by analytic constructions or continuity arguments,
separately from rounded tie stops. A stopped witness is an evaluated feasible
point, not an accepted estimate. Its likelihood can have finite bounds even
when the unclosed global upper bound remains infinite.

Work is controlled by $K_u\min(\mathtt{iterlim},100000)$, using 64-bit counters:
$K_u=1$ for a solver iteration, $128$ for a certificate segment, box or profile
point, and $100000$ for an exact enumeration cell. Inner solver loops have
their own solver caps. Exhausted certificate or enumeration work produces a
budget stop with available witnesses; solver exhaustion produces a numerical
nonconvergence outcome. Closed forms and special-function evaluators have
separate documented evaluation conventions. A larger budget can resolve an
unfinished cover but cannot resolve an unrepresentable coefficient or an
arithmetic comparison below available precision. These counts limit work
units, not uniform wall-clock cost: an evaluation may scan the whole sample.

\subsection{Evidence that should accompany a result}
We use ``proved'' for a real-arithmetic argument covering the stated sample
class; ``certified'' for a concrete sample covered by analytic inequalities
and computed comparisons; and ``numerical'' for exploratory roots, profiles
and simulations without an exhaustive exclusion argument. A proof sketch
below names the reduction and the inequality doing the work. The linked
project reports retain the longer derivations and computations.

The historical wrapper rule allowed a failed optimizer only after a documented
score gate, commonly total score norm below $10^{-3}$ in fitter coordinates.
That gate checks local stationarity. It supplies neither global optimality
nor the parameter-accuracy enclosure demanded here. Partially certified
families still contain numerical routes, so the contract is a design standard
and a scoped achieved guarantee, not a claim that every current successful
route already meets it.
\source{\code{dev/design/work-budgets.md}; \code{dev/validation/diagnostics.md}; family dossiers and the coverage inventory in Section~\ref{sec:software}.}

\section{Existence and uniqueness from the likelihood}\label{sec:exist}
Throughout this section a count sample has $x_i\in\{0,1,\ldots\}$,
$m=n^{-1}\sum_i x_i$ and $v=n^{-1}\sum_i(x_i-m)^2$. The divisor in $v$ is
$n$, not $n-1$. All-zero samples are separated whenever a zero-law boundary
makes coefficients unidentified.

\subsection{Negative binomial: a complete proof dependency}
Use the success-probability convention
\[
 f(x;r,p)=\frac{\Gamma(r+x)}{\Gamma(r)x!}p^r(1-p)^x,
 \qquad r>0,\quad 0<p<1.
\]
\begin{theorem}[Ordinary NB]\label{thm:nb}
For $m>0$, a finite joint NB maximum exists if and only if $v>m$, and is then
unique. For $v\le m$ the profile increases strictly to an unattained Poisson
supremum. All-zero samples instead attain a zero-law member in the package's
extended domain and have nonunique coefficients.
\end{theorem}
\begin{proof}[Proof sketch expanding Br\o ns and Jensen's convexity argument]
Strict conditional concavity gives $p(r)=r/(r+m)$. For $f=-\ell_{\rm prof}$,
\[
 f'=n\log(1+m/r)-\sum_i\sum_{j=0}^{x_i-1}\frac1{r+j},\qquad
 f''=-\frac{nm}{r(r+m)}+\sum_i\sum_{j=0}^{x_i-1}\frac1{(r+j)^2}.
\]
Set $b=r\log(1+m/r)/m$. It increases from zero to one and
$b''=-m/[r(r+m)^2]<0$. The finite-sum lemma in the dossier proves strict
convexity of
\[
 T_q(r)=r(r+m)^2\sum_{j=0}^{q-1}(r+j)^{-2},\qquad q\ge1.
\]
The proof splits the telescoping sum at $\min(q-1,\lfloor m\rfloor)$;
positive second derivatives of the remaining factors give strictness,
including $m<1$ and $q=1$. Hence
\[
 h=\frac{f''}{b''}=n(r+m)-\frac1m\sum_iT_{x_i}(r)
\]
is strictly concave and tends to $n(v-m)/m$. A strictly concave function
with a finite limit at infinity has $h'>0$: a nonpositive derivative would,
by strict decrease of the derivative, force a negative eventual slope.

Let $G=f'/b'$. Since $f'$ and $b'$ vanish at infinity, $G$ is the weighted
tail average of $h$ with positive weight $-b''$. Thus $G>h$ and
$G'=(b''/b')(h-G)>0$. Its endpoint limits are
\[
 G(0+)=-\infty,\qquad G(\infty)=n(v-m)/m.
\]
There is exactly one zero if $v>m$, and no zero otherwise. At equality $G$
is strictly below its zero limit at every finite $r$. Reversing the sign from
$f$ to $\ell$ proves the stated monotonicity and globality. Along $p(r)$,
the masses converge to Poisson masses with mean $m$.
\end{proof}
This is a classical conclusion \cite{anscombe,levin,aragon}, not a priority
claim. The project's proof report inspected the complete 1988 Br\o ns--Jensen
preprint \cite{brons}. Wang's abstract \cite{wang} reports a major problem with the 1992
proof. Independently, the accessible 1991 preprint has a fixed-coordinate
boundary cancellation gap: in its notation, at $M=3$ and held $a=F_2$,
\[
 u(z,0,a)=-\frac{a}{2(1+z)}.
\]
For $a=1/4$, the boundary values at $z=1,2$ are $-1/16,-1/24$, with difference
$1/48$, rather than zero. Replacing them by boundary values at $a=0$ changes
the integration slice. This diagnoses that preprint step; it is not a
counterexample to Theorem~\ref{thm:nb}.
\todo{Obtain Wang's full note and the published Arag\'on--Eberly--Eberly proof before identifying this preprint gap with Wang's exact objection.}
\source{\code{2026-11-07-nb-proof-report.md} and the NB dossier, ``Existence and uniqueness proof''. All dated report names here and below are relative to \code{dev/research/evaluations/}.}

\subsection{Poisson--inverse Gaussian}
For the PIG family let $\mu>0$ be the mean and $\lambda>0$ the mixing shape
in the dossier's parametrization. The Poisson face is $\lambda\to\infty$.
\begin{theorem}[PIG dispersion criterion]\label{thm:pig}
For a count sample with $m>0$, the free-mean PIG likelihood has a finite
maximum exactly when $v>m$, and that maximum is unique with $\widehat\mu=m$.
With $\mu$ held positive, the corresponding criterion is
\[
 v+(m-\mu)^2>m,
\]
and the finite shape maximum is again unique. Equality and the reversed
inequality give the Poisson supremum with the same held mean.
\end{theorem}
\begin{proof}[Proof sketch]
The finite half-integer Bessel polynomial $Q_{k-1}(1/z)$ has positive term
weights. Under their normalized distribution let $E_k,V_k$ be the mean and
variance of the term index. Its polynomial differential equation gives, for
$k\ge2$ and $c=k(k-1)$,
\[
 \frac{dE_k}{d\log z}=-V_k,\qquad
 V_k=c-(2z+1)E_k-E_k^2.
\]
The essential inequality is
\begin{equation}\label{eq:bessel}
 V_k< B(E_k):=\frac{E_k(k^2-E_k^2)}{k^2+E_k^2}.
\end{equation}
To verify it, put $F=V_k-B(E_k)$. At any zero,
\[
 zF'(z)=\frac{2E_k k(E_k+k)^2
 (E_k^3+E_k^2k-E_kk^2+k^3)}{(E_k^2+k^2)^3}>0.
\]
Here $0<E_k<k-1<k$, making the last polynomial positive. But the finite
polynomial expansion at infinity gives $F=-c/(2z^2)+O(z^{-3})<0$ eventually.
Every zero would cross upwards; returning to this negative tail is
impossible. This proves \eqref{eq:bessel}.

For held $\mu$ set
\[
 w=\frac{\sqrt{1+2\mu^2/\lambda}-1}{\sqrt{1+2\mu^2/\lambda}+1},
 \qquad z=\frac{\mu(1-w^2)}{2w},\quad 0<w<1.
\]
The individual log mass for $k\ge1$ is
\[
 -\mu(1-w)+k\log\frac{\mu(1-w)}{1+w}
 +\log Q_{k-1}(1/z)-\log(k!).
\]
Its second derivative, multiplied by $(1-w^2)^2$, is
\[
 -4kw+2E_k(3+w^2)+(V_k-E_k)(1+w^2)^2/w^2.
\]
Substitute \eqref{eq:bessel} and put $a=E_k/k$. Dividing by $k$, the upper
bound factors as
\[
 \frac{2(a-w)^2(aw^2-a-2w)}{w^2(1+a^2)}\le0.
\]
Strictness holds before this last weak inequality. For $k=1$ the scaled
second derivative is $-4w<0$; for $k=0$ the log mass is linear. Thus a sample
containing a positive count has a strictly concave $w$ likelihood. Its
one-sided derivative at zero is
\[
 n\mu-2\sum_i x_i+\frac1\mu\sum_i x_i(x_i-1)
 =\frac n\mu\{v+(m-\mu)^2-m\}.
\]
The likelihood tends to $-\infty$ at $w=1$. The sign at zero therefore
decides existence, including equality.

For the joint problem, the fixed-$z$ mean reduction is separate from this
held-mean calculation. With $t=\lambda/\mu$ and
$\lambda=(z^2-t^2)/2$, the $t$-dependent log likelihood is
$nt+(\sum_i x_i)\log\lambda$. Its unique conditional maximum gives $\mu=m$;
the inward derivative also excludes the infinite-mean face. The remaining
profile is precisely the held-$m$ problem just proved.
\end{proof}
The proof uses a stronger variance bound than $V_k\le E_k$. The latter is
useful for a held-positive-shape mean solve but does not establish joint
profile uniqueness. Decreasing terms in a score difference cannot by
themselves determine its crossings. The source report records that rejected
shortcut and the subsequent analytic repair.

The criterion and polynomial inequality were not located after the bounded
search. Willmot \cite{willmot} remains contextual metadata; failed retrieval
of Carlson is not evidence that the result is absent there.
\source{\code{2026-11-06-pig-cert-report.md}; its data directory's \code{concavity-proof.md}; PIG dossier.}

\subsection{Two concave shape profiles}
\begin{theorem}[Kumaraswamy]\label{thm:kumar}
For every nonconstant sample in $(0,1)$, the two-free-shape Kumaraswamy
likelihood has a unique finite global maximizer. Constant interior samples
have unbounded likelihood when both shapes are free.
\end{theorem}
\begin{proof}
Write $t_i=-\log x_i>0$, $T=\sum_i t_i$ and
$G(a)=\sum_i-\log(1-e^{-at_i})$. Profiling the second shape gives
$\widehat b=n/G(a)$ and
\[
 h(a)=n\log a-n\log G(a)-aT+G(a)+C.
\]
Each summand of $G$ is $\sum_{j\ge1}e^{-jat_i}/j$, so $G$ is a positive
Laplace transform. Cauchy--Schwarz gives $(\log G)''\ge0$. Moreover
\[
 G''(a)=\sum_i\frac{t_i^2e^{at_i}}{(e^{at_i}-1)^2}<n/a^2,
\]
because $2\sinh(u/2)>u$ for $u>0$. Hence
$h''=-n/a^2-n(\log G)''+G''<0$.
At zero $h(a)\to-\infty$. At infinity the leading linear term is
$a(nt_{\min}-T)$, strictly negative for nonconstant data. Both tails are
therefore coercive and there is one maximum. For constant data this linear
term vanishes and the remaining $n\log a$ divergence is unbounded.
\end{proof}
Jones \cite{jones} supplies the profiled likelihood; strict concavity and the
all-nonconstant-sample conclusion were not located in the inspected preview.
\source{Kumaraswamy dossier, ``Scalar certificate and uniqueness proof''.}

\begin{theorem}[Burr XII, scale held]\label{thm:burr}
Hold $s>0$ in the Burr XII density with shapes $c,k>0$. A nonconstant positive
sample containing an observation below $s$ has a unique finite two-shape
maximum. Constant samples at or below $s$ have unbounded likelihood. If all
observations are at least $s$ and some exceed it, the supremum is a Pareto
path at endpoint $s$, with the endpoint penalty specified below.
\end{theorem}
\begin{proof}[Proof sketch]
Let $z_i=\log(x_i/s)$,
$G(c)=\sum_i\log(1+e^{cz_i})$ and $R(c)=\sum_i\log(1+e^{-cz_i})$.
Then $\widehat k=n/G(c)$ and
\[
 \ell_p(c)=n\log n-n-\sum_i\log x_i-n\log\{G(c)/c\}-R(c).
\]
For $A(u)=\log(1+e^{-u})$, $u>0$, the dossier proves
\[
 (\log(A(u)/u))''\ge u^{-2}-e^{-u}>0.
\]
For positive log arguments add the constant term one to $A(u)/u$; sums
preserve log convexity. A zero argument contributes $\log2/c$.
Thus $G(c)/c$ is strictly log convex, and $R$ is convex, giving
$\ell_p''<0$. At $c\downarrow0$ the profile tends to $-\infty$.
If some $z_i<0$ and the sample is nonconstant, its large-$c$ asymptotics also
tend to $-\infty$; a negative log observation supplies the linear decay when
positive logs occur, while an all-negative nonconstant sample has
$\sum z_i<n\max z_i$. Constant nonpositive logs instead give divergence.

When every $z_i\ge0$, let $S=\sum_i z_i>0$ and
$Z=\#\{i:x_i=s\}$. The large-$c$ limit is
\[
 n\log(n/S)-n-\sum_i\log x_i-Z\log2.
\]
Strict concavity and this finite limit imply that the profile increases
strictly to it. Each tied endpoint observation contributes the extra
$-\log2$, so the reachable likelihood is not simply the closed Pareto
likelihood.
\end{proof}
The exact overlap with Wingo \cite{wingo83,wingo93} is open because only the
abstracts were inspected. No priority claim for Burr uniqueness is warranted.
Reciprocal and threshold transformations transport the result only with the
held scale and the likelihood Jacobian correctly retained.
\source{\code{2026-11-06-burr-scale-report.md} and Burr dossier.}

\subsection{Fr\'echet and minimum multiplicity}
Use the dossier's Fr\'echet density with positive held shape $\lambda$,
location $\mu<\min_i x_i$ and scale $\sigma>0$:
\[
 f(x)=\lambda\sigma^{\lambda}(x-\mu)^{-\lambda-1}
       \exp\{-[\sigma/(x-\mu)]^\lambda\}.
\]
Here $\sigma$ is the public ordinary scale, so $(X-\mu)/\sigma$ has
unit-scale Fr\'echet distribution.
\begin{theorem}[Held-shape Fr\'echet]\label{thm:frechet}
Let $k$ be the multiplicity of the sample minimum. With location and scale
free and $\lambda$ held, a finite MLE exists if and only if
$\lambda(n-k)>k$, and is then unique. At equality the supremum is finite
but unattained; below equality the likelihood is unbounded.
\end{theorem}
\begin{proof}
Set $\epsilon=\min_i x_i-\mu$, $g_i=x_i-\min_j x_j$ and
$S_p=\sum_i(g_i+\epsilon)^{-p}$. Profiling gives
$\widehat\sigma^\lambda=n/S_\lambda$. The gap score has the sign of
\[
 n\lambda R-(\lambda+1),\qquad
 R=\frac{S_{\lambda+1}}{S_\lambda S_1}.
\]
For $q_i=(g_i+\epsilon)^{-1}$ put
$E_p=\sum_iq_i^{p+1}/\sum_iq_i^p$. Weighted covariance of $q$ and $\log q$
is strictly positive for nonconstant data, so $E_p$ increases with $p$.
Consequently
\[
 \frac{d\log R}{d\epsilon}
 =-(\lambda+1)E_{\lambda+1}+\lambda E_\lambda+E_1<0.
\]
The limits of $R$ are $1/k$ and $1/n$. A unique downward score crossing
exists precisely under the stated strict inequality. At the small gap,
the coefficient of $\log\epsilon$ in the profile is
$\lambda(n-k)-k$. This gives unboundedness below the threshold and a finite
limit at equality; the strictly negative gap score at equality excludes
attainment. Constants have $k=n$ and fall in the unbounded case.
\end{proof}
For equality, the finite supremum is
\[
 n\log(n\lambda)-n-n\log k
   -(\lambda+1)\sum_{g_i>0}\log g_i.
\]
At held $\lambda=1$, $(0,0,1,2)$ gives $2\log2-4$, unattained;
$(0,0,0,1)$ is unbounded. Holding scale as well gives a unique finite
location fit even for constants: the relevant ratio
$S_{\lambda+1}/S_1$ decreases from infinity to zero.

The searched Fr\'echet theorem with location fixed \cite{frechet2016} does not imply this
minimum-multiplicity criterion. The exact threshold and equality case were
not located after search; full comparisons with Smith \cite{smith}, Cheng
and Iles \cite{cheng}, and Hall and Wang \cite{hallwang} remain open.
\source{\code{2026-11-07-evt-scalar-report.md} and Fr\'echet dossier.}

\subsection{Conditional count likelihoods}
\begin{theorem}[ZTNB and conditional ZINB]\label{thm:conditional}
For zero-truncated NB with $0<p<1$ held, the size log likelihood is strictly
concave on $r>0$; on the extended domain $r\ge0$ the maximum is unique for
positive samples of mean greater than one, possibly at the logarithmic-series
member $r=0$. For ZINB data containing a positive count, free inflation with
either positive size held or interior probability held has a unique maximizer.
\end{theorem}
\begin{proof}
With $t=-\log p$, the ZTNB size curvature is
\[
 -\frac n{r^2}-\sum_i\sum_{j=1}^{x_i-1}(r+j)^{-2}
       +\frac{nt^2}{4\sinh^2(rt/2)}<0.
\]
The inequality again follows from $2\sinh(u/2)>u$. Continuity at zero and
vanishing likelihood at infinite size complete the extended-domain claim.
All-one samples have the separate point-mass nonidentification when
probability is free.

For ZINB put $z=n_0/n$. At fixed base parameters the inflation maximum is
\[
 \widehat\pi=\max\{0,(z-p^r)/(1-p^r)\}.
\]
In the positive-inflation region the likelihood separates into a saturated
zero/nonzero likelihood and a positive-count truncated-NB likelihood.
With size held, ordinary and truncated NB are strictly concave in the
probability natural parameter; their scores agree at $p^r=z$, giving a
unique joined maximum. With probability held, both size branches are
strictly concave by the displayed curvature and the ordinary NB curvature.
Their scores again agree at the join. Positive observations make the ZINB
likelihood vanish at either size extreme in this held-probability problem.
This gives existence and uniqueness.
\end{proof}
The separation is partly known \cite{haslett}; the held-parameter uniqueness
theorems were not located after search. Separate concavity does not establish
joint size/probability uniqueness. The latter distinction leads to
Section~\ref{sec:zinb}.
\source{ZTNB and ZINB dossiers; \code{2026-11-06-zinb-cert-report.md}.}

\section{Explicit failures of uniqueness}\label{sec:modes}
A multiple-root likelihood, a tie between distinct distributions and
coefficient nonidentification are different phenomena. We state each example
with its evidence level. Decimal values illustrate the retained calculations;
only exact definitions or the stated conditional certificates support exact
claims.

\subsection{An exact generalized Pareto tie}
For strictly positive excesses, threshold zero and $\xi\ge-1$, use
$f(y;\sigma,\xi)=\sigma^{-1}(1+\xi y/\sigma)^{-1-1/\xi}$, with the
exponential limit at zero and the closed uniform endpoint at $\xi=-1$.
Grimshaw's scalar reduction is established prior work \cite{grimshaw}.
\begin{theorem}[GPD scalar reduction]\label{thm:gpdprofile}
For $b>-1$, write $M=\max_i y_i$, $v_i=y_i/M$, $b=\xi M/\sigma$, and define
\[
 D(b)=\overline{\int_0^{v_i}(1+bs)^{-1}\,ds}.
\]
It is continuous through $b=0$. The conditional joint optimum is
$\xi=bD$, $\sigma=MD$ when $bD\ge-1$, with mean log likelihood
$-\log M-\log D-bD-1$. Otherwise the constrained uniform branch cannot
beat the uniform member on $[0,M]$.
\end{theorem}
\begin{proof}
At fixed $b$ put $u=\sigma/M$. The mean log likelihood is
$-\log M-\log u-bD-D/u$. Its derivative in $u$ is $(D-u)/u^2$,
so the unique unconstrained maximum is $u=D$. If it violates
$\xi=bu\ge-1$, the constrained maximum is the uniform branch
$\xi=-1$, $\sigma=-M/b\ge M$. Its density cannot exceed $1/M$.
The integral representation defines $D$ and this optimization continuously
at $b=0$.
\end{proof}

\begin{proposition}[GPD equal global maxima]\label{prop:gpd}
Put $b(k)=ke^{1+k}$ and
\[
 a(k)=\frac{\exp\{[3k-\log(1+b(k))]/2\}-1}{b(k)},\qquad
 a(0)=(3/e-1)/2.
\]
Let $a_* =\max_{k>-1}a(k)$ and choose any maximizer $k_*$. For the sample
$(a_*,a_*,1)$ the distinct members $(\sigma,\xi)=(1,-1)$ and
$(e^{-1-k_*},k_*)$ are global maximizers, each with total log likelihood zero.
\end{proposition}
\begin{proof}[Proof sketch]
The continuous function $a$ tends to $-\infty$ at $-1$ and to zero at
infinity, is less than one everywhere, and has
$a'(0)=3(e^2-2e-1)/(8e)>0$. Thus $0<a_*<1$ is attained away from zero.
At fixed $b$, $h=-\log D-bD-1$ decreases with $D$ on its admissible
range, while $D$ increases with the smaller datum $a$ with derivative
$2/[3(1+ba)]$. Solving $h=0$ gives $D=e^{-1-k}$ and $b=ke^{1+k}$;
substitution gives precisely $a=a(k)$. Since $k\mapsto b(k)$ maps
$(-1,\infty)$ increasingly onto $(-1,\infty)$, $a_*\ge a(k)$ bounds
every admissible profile value by zero. The constrained uniform branch is
also bounded by zero. Both displayed members attain the bound.
\end{proof}
The numerical illustration is
\[
\begin{split}
 a&\approx0.07880675592383729,\\
 \xi&\approx0.5808198884419236,\qquad
 \sigma\approx0.2058062908136120.
\end{split}
\]
The max-defined sample, not its rounded approximation, defines the exact tie.
Multiple GPD profile roots are known; this restricted-endpoint global tie was
not located in the inspected Grimshaw text.
\source{GPD dossier, ``Nonuniqueness''; \code{2026-11-05-gpd-cert-data/check_uniqueness.py}.}

\subsection{Zeros supporting two explanations}
\begin{proposition}[ZINB held-inflation crossover]\label{prop:zinbtie}
For 99 zeros and one count of ten, with size $r=10$ held, the probability
likelihood can have two local maxima. Their order reverses as held inflation
passes a value near $0.7683947756875808$. A value in the certified crossover
interval has two equal global maxima.
\end{proposition}
\begin{proof}[Branch count and continuity]
For general $n_0$, $k=n-n_0>0$, $S=\sum_i x_i$, let
$a=\pi/(1-\pi)$, $u=kr/(kr+S)$ and $v=nr/(nr+S)$. The probability score
has the sign of
\[
 a[kr-(kr+S)p]+p^r[nr-(nr+S)p].
\]
It is positive below $u$ and negative above $v$. Between them its zeros
solve
\[
 R(p)=\frac{p^r(nr+S)(v-p)}{(kr+S)(p-u)}=a.
\]
The sign of $R'$ is that of
$-rp^2+[r(u+v)-(v-u)]p-ruv$. Its two possible turning points partition
all monotone branches, so at most three score roots occur. The retained
branch comparisons give opposite orderings of the two maxima at inflation
values bracketing the crossover. Endpoints have zero likelihood, and all
branches are covered. Continuity gives a global equal-height value between
the brackets. Rounded agreement at the displayed inflation is not itself
the exact tie proof.
\end{proof}
Structural versus sampling zeros are familiar \cite{ridout98}. The fixed
inflation three-root bound and equal-height crossover were not located after
search. Freeing inflation instead leads to Theorem~\ref{thm:conditional}.
\source{ZINB dossier; \code{2026-11-06-zinb-cert-report.md}; \code{2026-11-06-nbfam-cert-data/zinb-tie.json}.}

\subsection{An analytic SGED scale tie}
\begin{proposition}[Two global SGED scales]\label{prop:sged}
Take any number of repeated observations $x_i=-1$, held mean zero and
held power $\nu=1/2$ in the standardized Fern\'andez--Steel SGED. Let
$R\in(4,6)$ solve $\log R=2(1-1/R)$, put
\[
 K=2\sqrt{R-1}/R,\qquad \xi=\sqrt{1+K^2/6}.
\]
At this held skew, two separated scales are global maximizers, although no
observation equals the held mean.
\end{proposition}
\begin{proof}
Use the dossier's dimensionless inverse-scale coordinate $r$, in which the
observation crosses the kernel cusp at $r=1$. Beyond the cusp, the cost
relative to that point is
\[
 n\{-\log r+K\sqrt{r-1}\},\qquad r\ge1.
\]
The score equation is quadratic in $\sqrt{r-1}$; the first stationary
point is a maximum of cost, the second a minimum at $R$. The defining
equation makes the cost at $R$ equal the zero cost at the cusp. Cost tends
to infinity on the far tail and decreases towards the cusp on $(0,1)$.
Thus no lower cost occurs, and the two scales tie globally.
\end{proof}
The retained 70-digit check gives $R=4.921553634567505$ and
$\xi=1.0525849322732358$ as numerical illustrations. The SGED scale map
uses fixed mean and standard deviation coordinates, not a held raw mode.
Schoups and Vrugt \cite{schoups} give related standardization formulas;
their stated power range excludes this subunit example. This tie was not
located after search.
\source{\code{2026-11-07-sged-subunit-report.md} and SGED dossier.}

\subsection{A beta-prime tie obtained by continuity}
For shape $a$ held, free shape $b$ and scale $s$, the beta-prime density is
$x^{a-1}s^{-a}/\{B(a,b)(1+x/s)^{a+b}\}$.
\begin{proposition}[Conditional beta-prime global tie]\label{prop:betapr}
Hold $a=0.22690393135372058$ and take the nine observations
\begin{align*}
(&1.9901008294052028\,10^{-5},\ 2.209857117039297\,10^{-5},\\
 &0.0002450722434289249,\ 12.288342947780086,\\
 &942.976056905759,\ 16.239512580216314,\\
 &2.898690421966857,\ 1.2684036923492177,\ 400970.63293690106).
\end{align*}
Multiply the last observation by $e^t$. Conditional on the retained endpoint
certificates and uniform gap comparisons, some $t\in(5.6463,5.6466)$ gives
at least two separated global maxima in $(b,s)$, on opposite sides of
$b=0.12$.
\end{proposition}
\begin{proof}[Certificate and continuity]
For each $b$, the log-scale score has a unique root. Profile derivative and
Schur-complement bounds isolate competing maxima without assuming profile
unimodality. The endpoint certificates select opposite sides of $b=0.12$.
Throughout the $t$ bracket, Lipschitz bounds keep a finite witness at least
$0.17569$ above the barrier profile and $17.33799$ above the gamma face;
both shape exteriors are uniformly excluded. On the resulting compact sets
the two side maxima are continuous in $t$. Their difference changes sign,
so it vanishes somewhere. The barrier and face gaps keep the tied maxima
separated and global.
\end{proof}
At $t=0$ the two local profiles are approximately
\[
\begin{array}{rrr}
 b&s&\ell\\
 0.09175915065&0.000500743658&-29.14985912088\\
 0.22150182273&2.54933710323&-28.55611432402.
\end{array}
\]
Near the numerical crossing $t=5.64643839503$ the implementation gives a tie
stop with two witnesses. That decimal is not asserted to be an exact
crossing. Tulupyev et al.\ \cite{tulupyev} hold the numerator shape at 2;
their critical-point analysis does not settle this arbitrary-held-shape
example. An equivalent global tie was not located after search.
\source{\code{2026-11-07-betapr-prof-report.md}; \code{counterexample-refined.json} in its data directory.}

\subsection{Beta negative binomial: labeled nonuniqueness}
The pmf
\[
 f(x;r,\alpha,\beta)=\frac{(r)_x}{x!}
             \frac{B(\alpha+r,\beta+x)}{B(\alpha,\beta)}
\]
is invariant under exchanging $r$ and $\beta$. This known nonidentifiability
\cite{ueki} alone does not exclude a maximum on the fixed diagonal.
\begin{proposition}[An off-diagonal BNB maximum]\label{prop:bnb}
Hold $\alpha=2$ and use the sample specified by the following count-frequency
pairs (all omitted counts have frequency zero):
\[
\begin{array}{rrrrrrrrr}
 x&0&1&2&3&4&5&6&7\\
 n_x&135&84&57&39&35&38&18&14\\[2pt]
 x&8&9&10&11&12&13&14&15\\
 n_x&14&7&4&9&6&6&3&1\\[2pt]
 x&16&17&18&20&21&22&23&28\\
 n_x&3&2&2&2&1&3&1&1\\[2pt]
 x&29&31&35&36&42&46&47&48\\
 n_x&1&1&2&2&1&1&1&1\\[2pt]
 x&51&64&125&175&&&&\\
 n_x&1&2&1&1&&&&
\end{array}
\]
These 500 observations have at least two distinct labeled global maximizers
in $(r,\beta)$, representing the same law, conditional on the retained
diagonal-exclusion comparisons.
\end{proposition}
\begin{proof}[Proof sketch]
The fixed-$\alpha$ likelihood attains its maximum: a shape tending to zero
makes every positive-count mass vanish uniformly in the other shape; a
shape diverging with the other bounded away from zero gives stochastic
escape. The report's superlevel-set argument excludes the simultaneous
escapes too. The witness $(r,\alpha,\beta)=(1,2,5)$ has
$f(x)=60/\{(x+5)(x+6)(x+7)\}$ and total log likelihood
$-1276.345658177647696$. The upper bound over the entire diagonal is
$-1276.445827365689238$, leaving a gap at least $0.10016918804154$.

For diagonal shape $s$, factor the mass as $f(0;s)R_x(s)$; the first
factor decreases, the second increases. Thus
$f(x;s)\le f(0;l)R_x(h)$ on $[l,h]$. The 60-digit certificate covers 861
segments of $[10^{-6},256]$ and bounds both exterior tails, with absolute
log-likelihood padding $10^{-20}$. No diagonal maximizer remains. An
attained off-diagonal maximum and its swap are distinct labeled maximizers.
The witness itself is not claimed to maximize.
\end{proof}
Scalar uniqueness with only one free shape is proved separately by a
single-crossing digamma-integral argument. It cannot be substituted for
multivariate uniqueness. Uniqueness after imposing $r\le\beta$ remains open.
\source{\code{2026-11-06-bnbinom-cert-report.md}; its \code{nonuniqueness.json}; the frozen BNB \code{golden.json}.}

\subsection{Multimodality without equal global values}
\begin{proposition}[One-observation paralogistic likelihood]\label{prop:para}
With scale held at one and $x=e^{-0.001}$, the paralogistic shape likelihood
has stationary points near
\[
\begin{array}{rrl}
 a&\ell&\text{type}\\
 2.895379787257729&-0.5700990971060364&\text{local maximum}\\
 1140.924271780363&-303.6829139605954&\text{local minimum}\\
 9262.778134534774&8.126729864909167&\text{global maximum}.
\end{array}
\]
The globality statement is conditional on the certificate's computed
comparisons. The root values have an independent 70-digit numerical check using the
stored binary64 evaluation of $e^{-0.001}$.
\end{proposition}
The held-scale log density, with $z=\log x$, is
\[
 2\log a+(a-1)z-(a+1)\log(1+e^{az}).
\]
Differentiation supplies the score and curvature. The certificate combines
root inclusion, likelihood bounds and analytic tail exclusion; the three
numerically stationary points alone would not prove completeness. A single
observation can therefore invalidate an assumed one-root shape solve.
Starting-values documentation gives the likelihood and gradients, but this
one-observation result was not located after search.
\source{\code{2026-11-06-paralogis-modes-report.md}; its \code{oracle.py} and \code{modes.tsv}.}

\begin{proposition}[Location-scale $t$ profiles]\label{prop:t}
For the following weighted samples, a location-scale $t$ likelihood with
ordinary scale and $\nu\ge1$ has the indicated certified global interior
maximum, conditional on the retained comparisons:
\[
\begin{array}{rr|rr}
\multicolumn{2}{c|}{\text{two interior maxima}}&
\multicolumn{2}{c}{\text{stable Cauchy competitor}}\\
 x&n_x&x&n_x\\
 -2.5267827738984026&2&-3.3753182637169354&1\\
 -1.3086958612269082&1&-1.2713352565886626&11\\
 -0.917311177389449&4&-0.8001089566898665&8\\
 -0.7470612434645413&4&0.4550011648748676&11\\
 -0.6822670978323393&1&0.4631607376683885&7\\
 0.5175256317384838&8&0.48440591294209223&8\\
 0.8692557396694677&11&1.2931658287635732&9
\end{array}
\]
The first has two interior profile maxima, near $\nu=1.0102899103927294$
and $49.705241724$; the second has its global maximum at $33.905406907$,
although the profile initially decreases from the restricted Cauchy member.
\end{proposition}
At fixed shape, certified nuisance solutions, profile envelopes and exterior
bounds compare all shapes and the normal and lower faces. The final replay
certifies the larger-shape winner in both cases. For the first sample the
small maximum has log likelihood $-46.21889676163025$, while the larger has
$-43.980462713143716$ in the original numerical fixture. For the second,
the Cauchy endpoint has log likelihood $-83.65806430406303$ and the interior
maximum $-77.96128211222002$. The final certificate report supplies the
winning shapes above; exploratory root decimals do not assert exact roots.

The standardized-$t$ domain $\nu>2$ uses the same ordinary-scale likelihood
with a different coefficient map; both retained winners are also certified
there. General profile unimodality is not assumed. Fern\'andez and Steel
\cite{fernandez} supply classical context for $t$ boundary pathologies;
the inspected abstract does not establish these interior examples.
\source{\code{2026-11-06-t-close-report.md}; \code{2026-11-04-t-closure-data/fixtures.json}.}

For comparison, the half-$t$ sample
$(0.12773,12.5979,3.91782,2.58658,0.138212)$ has numerical profile maxima
at $\nu=0.41508$ and $1.17074$, with log likelihoods $-11.72979$ and
$-11.98253$. A heterogeneous scan found 32 multimodal profiles among 1200
samples. This is a numerical finding, not a prevalence estimate under a
single sampling design. The failure gallery's rounded sample and scan do
not supply a universal half-$t$ mode-count theorem. For completeness, with held mean zero, $\nu=3$ and $\xi=10$, the six-point
Fern\'andez--Steel skew-$t$ sample
\begin{align*}
 (&-35.321072334324725,\ 11.459074428000962,\\
 &-0.015351125824354262,\ -0.2513099360462538,\\
 &-297.3459307209368,\ 0.5881721220894813)
\end{align*}
has the following numerical stationary points:
\[
\begin{array}{rrl}
 \log\sigma&\ell&\text{type}\\
 3.764665559514996&-50.43173618293705&\text{maximum}\\
 4.850139380730862&-51.428090712695514&\text{minimum}\\
 5.894765235143508&-38.8073025560888&\text{maximum}.
\end{array}
\] At symmetry the scale score
is strictly monotone; skewing with the public mean held creates the extra
mode. This numerical construction is not an exact-tie theorem.
\source{\code{2026-11-06-sstd-cert-report.md}; its \code{skew-scale.json}; \code{failure-gallery.tex}, half-$t$ section.}

\section{Boundaries, reachable likelihoods and limit models}\label{sec:zinb}
\subsection{A face score can be sufficient without being decisive}
Let $n_0$ be the zero count, $k=n-n_0>0$, $S=\sum_i x_i$,
$z=n_0/n$, $q=k/n$ and $m_+=S/k$. In this section the full-sample mean is
$\bar x=S/n$, so $m_+$ is not the $m$ of Theorem~\ref{thm:nb}.
Profiling inflation yields either the ordinary NB branch or
\[
 B+T(r),\qquad B=n_0\log z+k\log q,
\]
where $T$ is the probability-profiled positive-count ZTNB log likelihood.
After maximizing the positive-count probability, $B+T(r)$ is an upper
bound for every fixed-size ZINB fit, including those for which saturation
would imply negative inflation.

\begin{theorem}[Established all-free ZINB regions]\label{thm:zinb}
For data containing a positive count:
\begin{enumerate}
\item If $z\le e^{-\bar x}$, a finite ZINB maximum exists exactly when
$v>\bar x$ and is then unique, at inflation zero.
\item Let $(\lambda_*,\pi_*)$ be the admissible maximized ZIP fit and put
$A=\sum_i x_i(x_i-1)/2$, $C=A-S\lambda_*/2$. Then $C>0$ is sufficient
for a finite ZINB maximum. On the positive-inflation ZIP face, $C\ne0$
for integer empirical data.
\item If every positive count equals $j\ge2$, define
$\lambda_j/(1-e^{-\lambda_j})=j$. If $q\le1-e^{-\lambda_j}$,
no finite ZINB maximum exists.
\end{enumerate}
\end{theorem}
\begin{proof}
The ordinary NB zero mass at its fixed-$r$ probability optimum is
$(1+\bar x/r)^{-r}>e^{-\bar x}$. In the zero-deficit region, inflation is
therefore zero for every $r$; Theorem~\ref{thm:nb} applies.

Along $p=\exp(-\lambda_*/r)$ with inflation $\pi_*$ held,
\begin{equation}\label{eq:zipgain}
 \ell_{\rm ZINB}-\ell_{\rm ZIP}=C/r+O(r^{-2}).
\end{equation}
For each positive count the first coefficient is
$x(x-1)/2-x\lambda_*/2$; the zero mass is exactly unchanged.
If $C>0$, a finite witness beats the global ZIP face. The size profile is
continuous, tends to that face at infinity and to $-\infty$ at zero, giving
attainment. At positive ZIP inflation, $C=0$ would make
$\lambda_*=2A/S$ rational and the conditional-mean equation would make
$e^{-\lambda_*}=1-\lambda_*/m_+$ rational. Hermite--Lindemann excludes
this for positive rational $\lambda_*$.

For the third assertion, an NB positive-count law is a nondegenerate
mixture of truncated Poisson laws. The conditional mass at $j$ has a
unique maximum at $\lambda_j$. Every nondegenerate mixture has strictly
smaller mass there. The assumption on $q$ makes the saturating ZIP fit
admissible; it reaches the universal indicator/conditional-mass bound in
the closure. No finite NB mixture attains it.
\end{proof}
Thus $(0,0,0,4)$ has $\bar x=1$, $v=3$, yet no finite ZINB maximum. The
ordinary NB variance rule cannot be transferred by analogy. The
admissibility requirement is also essential: with $(1,4)$ the unconstrained
positive-count Poisson fit suggests a positive inward coefficient, but
would require negative inflation. The admissible ZIP fit has
$\lambda_*=5/2$, $\pi_*=0$ and $C=-1/4$; the zero-deficit theorem correctly
gives nonexistence.

In the excess-zero region, the full profile has a single join from ordinary
NB to saturated ZTNB because $(1+\bar x/r)^{-r}$ decreases strictly from
one to $e^{-\bar x}$. The profile and its derivative agree at the join.
For $C<0$, \eqref{eq:zipgain} establishes only a local preference for ZIP
near infinite size. It does not exclude a distant finite maximum. The
published ZIP-versus-ZINB score-test connection is partly known
\cite{ridout01}; it does not by itself supply this global exclusion.
\source{\code{2026-11-07-zinb-theory-report.md}; ZINB dossier.}

\subsection{A certified finite optimum too close to the face}
\begin{proposition}[Near-ZIP sample]\label{prop:near}
For 5000 zeros, 3958 ones, 4063 twos and 1979 fours, the independent
65-digit padded certificate gives a unique finite global ZINB maximum:
\[
\begin{array}{lr}
 r&8108.2388581481329544584742965\\
 p&0.999803508015713486829597414531\\
 \pi&0.163276396309303523872445174295.
\end{array}
\]
The total log likelihood comparison is
\[
\begin{array}{lr}
 \text{finite ZINB}&-22998.0004858914916933762597635\\
 \text{ZIP face}&-22998.0005325865669794057284629\\
 \text{finite gain}&0.0000466950752860294686993812113.
\end{array}
\]
\end{proposition}
The certificate is conditional on comparisons padded by $10^{-45}$, not on
outward rounding. With dispersion $d=1/r$ and $u=-r\log p$, the positive
likelihood is evaluated without large gamma differences:
\begin{align*}
 L(d,u)={}&\sum_{x_i>0}\sum_{j=0}^{x_i-1}\log(1+jd)
 +S\log u+S h(ud)-ku\\
 &-k\log(1-e^{-u})-\sum_{x_i>0}\log(x_i!),\qquad
 h(v)=\log\{(1-e^{-v})/v\},\quad h(0)=0.
\end{align*}
A compact cover excludes sizes below 100. On the remaining region the
reported bounds prove a negative-definite Hessian in a dispersion/mean box,
and a signed bracket encloses the sole maximum. This is a global
sample-specific argument; it does not prove the general ZINB conjecture.

The historical fitter's face allowance $10^{-8}n=0.00015$ is larger than
the complete finite gain. That explains its false ZIP classification.
The certified fitter instead stops on budget when its cover cannot be
closed. At the optimum the profile curvature in $\log r$ is
$-9.33482447\times10^{-5}$; the total observed negative Hessian in
$(\log r,\log u,\pi)$ has condition number about $5.29\times10^8$.
These coordinate-dependent optimization diagnostics are not a claim about
Fisher covariance. The gain per observation is $3.11\times10^{-9}$.
Increasing work alone cannot repair a size-score comparison beneath its
practical binary64 resolution.
\source{\code{2026-11-07-zinb-theory-report.md}, Sections 4--5 and \code{near-certificate.json}.}

\subsection{Reachable boundary likelihood is path-dependent}
Within a common dominating measure, the Hellinger closure of the family
supplies candidate limiting models such as Poisson and exponential faces.
Concentration paths with unbounded likelihood require separate analysis;
they need not converge to a dominated law in that closure. Closure of
probability laws does not determine the pointwise density at an observed
support endpoint. Even weak
convergence, and even a density equivalence almost everywhere, leaves those
values undecided. A certificate must use the likelihood reachable through
the actual held-coordinate paths.

For standardized Fern\'andez--Steel skew $t$ with mean and standard deviation
held, shrink one half-width with degrees of freedom $\nu$ fixed. Define the kernel moment
$h_\nu=\Gamma((\nu-1)/2)/\{\sqrt\pi\Gamma(\nu/2)\}$, the first absolute
moment of $t_\nu/\sqrt\nu$, and put
\[
 v=1/(\nu-2)-h_\nu^2,\qquad
 \rho=\frac{h_\nu}{2(\nu-2)v},
\]
The same kernel has second moment $1/(\nu-2)$. Each observation at the resulting half-$t$ endpoint has a
multiplier $(1+\rho^2)^{-(\nu+1)/2}$ relative to its closed limiting density.
It is $0.6042259416$ at $\nu=3$. At the joint normal corner the largest
reachable multiplier is
\[
 e^{-\rho_N^2/2}=0.5473573954,\qquad
 \rho_N=\frac{\sqrt{2/\pi}}{2(1-2/\pi)}.
\]
Allowing degrees of freedom to vary toward a finite half-$t$ face can restore the full endpoint factor;
so can releasing appropriate moment holds. In Azzalini's skew $t$, free
location can approach the favorable endpoint side; location held exactly
there instead halves each tied observation's density. The distinction is
between paths satisfying constraints, not between equivalent labels for
the same unconstrained limiting law.
\todo{Choose whether to retain the endpoint-factor derivation in the paper or a supplement.}
\source{\code{2026-11-07-skew-theory-report.md}, especially its endpoint-path proofs.}

A second warning is supplied by standardized $t$. For masses $2/3$ at $b$
and $1/6$ at each of $b\pm d$, the normal-seam first coefficient is exactly
zero and the second coefficient is $n/4>0$. Nevertheless the optimized
profile decreases strictly with $\nu$ on $(2,\infty)$ and its supremum
is the collapse value $-n\log(2d)$ at $\nu\downarrow2$. Its ordinary
scale satisfies $s^2=d^2(\nu-2)/(3\nu)$ while standard deviation stays
$d/\sqrt3$. A zero boundary slope does not certify a face, and a finite
limiting standard deviation does not exclude concentration in the kernel
scale. The exact three-point result was not located after search.
\source{Standardized-$t$ dossier and \code{2026-11-05-t-root-data/derivation.md}.}

\section{Interior regularity and inference}\label{sec:inference}
A global MLE certificate is not an inferential calibration theorem. It does
not assert differentiability in quadratic mean (DQM), positive information,
consistency, an asymptotic distribution, or validity of a numerical Hessian.
The project's pointwise classification considers these questions separately
for 99 all-free family models; roughly 80 are regular and smooth, but that
summary is not a count of implemented covariance procedures.

Use $H^2(p,q)=\int(\sqrt p-\sqrt q)^2$, without a factor $1/2$. At a regular
point, per-observation information $I$ satisfies
\[
 H^2(p_\theta,p_{\theta+h})=h^TIh/4+o(\|h\|^2).
\]
Write the two sample matrices as
\[
 J=\sum_i s_i s_i^T,\qquad H=-\sum_i\partial^2\log f_\theta(x_i).
\]
Their distinct consistency
conditions matter. The outer product of the total score is not $J$.

\subsection{Symmetry and moving joins}
Azzalini skew normal has singular information at symmetry
\cite{hallin}. The mean/standard-deviation parametrized
Fern\'andez--Steel two-piece normal behaves differently. At standard
location and scale and $\xi=1$, in coordinates $(\mu,\log\sigma,\log\xi)$,
\[
 I=\operatorname{diag}(1,2,3-8/\pi).
\]
With $m_0=\sqrt{2/\pi}$, the skew score is
$\operatorname{sign}(z)z^2-2m_0z$. Its covariance with the location score
is $\E|Z|^3-2m_0\E Z^2=0$; symmetry makes it orthogonal to scale.
Its variance is $3-8/\pi>0$. Thus it is regular at symmetry. Nonsingular
two-piece information is partly known \cite{rubio}; the displayed constant
is a coordinate specialization, not a priority claim.

A moving join in a second derivative need not break DQM. The two-piece
normal, skew $t$ and Huber model have regular interior regimes with a valid
piecewise Hessian. But a generic finite-difference stencil crossing a join
need not have the formal smooth error order. A density theorem and a
validated production derivative method are separate dependencies.

\subsection{Cusp thresholds and endpoints}
For the exponential-power kernel
$f(x)=\nu\exp(-|(x-\mu)/a|^\nu)/\{2a\Gamma(1/\nu)\}$,
location information is finite exactly for $\nu>1/2$, with
\[
 I_{\mu\mu}=\frac{\nu^2}{a^2}
             \frac{\Gamma(2-1/\nu)}{\Gamma(1/\nu)}.
\]
The ordinary location Hessian has the correct integrable expectation only
for $\nu>1$. When $1<\nu\le2$ and location and shape are jointly fitted,
the source report leaves fitted-Hessian consistency open. At the critical
held power $\nu=1/2$, in its $a=1$ normalization, likelihood information at
rate $\sqrt{n\log n}$ is $1/16$, while raw $J/(n\log n)\to1/8$.
Rescaling the raw outer-product estimator would require its own argument.
The thresholds above $1/2$ are partly known from exponential-power work \cite{bottazzi};
the critical constants and fitted-Hessian issue were not located there.

For density proportional to $d^{\alpha-1}$ near a moving endpoint, with
controlled differentiated remainders, splitting the Hellinger integral at
distance of order $|h|$ gives
\[
 H^2\asymp
 \begin{cases}
 |h|^\alpha,&\alpha<2,\\
 h^2\log(1/|h|),&\alpha=2,\\
 h^2,&\alpha>2.
 \end{cases}
\]
These orders attach to endpoint-moving coordinates. At GEV shape $-3/4$,
one endpoint coordinate is fast at $n^{3/4}$; the two complementary smooth
coordinates retain $\sqrt n$. At the critical endpoint exponent two, the
endpoint experiment is Gaussian at $\sqrt{n\log n}$, not at $\sqrt n$.
Smith \cite{smith} supplies established nonregular context; exact coordinate
specializations and uniform transfer to each numerical estimator need
separate checking \cite{bucher}.

Other interior rank losses include BNB's $r=\beta$ and generalized beta's
second shape equal to one with its upper endpoint held. Tukey lambda's
uniform members at shapes one and two instead have a singular formal
pointwise score second moment; the all-free endpoint-moving model is not
DQM there. Formal scores must not be called regular Fisher information. The BNB symmetry explains a vanishing antisymmetric score on the
diagonal, but does not make the entire model regular after arbitrary
relabeling. Fixing a redundant parameter and estimating all of them are
different inferential problems.
\source{\code{2026-11-04-infer-moving-report.md}; \code{2026-11-04-dqm-a-report.md}, \code{2026-11-04-dqm-b-report.md}, \code{2026-11-04-dqm-c-report.md}; canonical pointwise inference classification.}

\subsection{Local experiments at closure points}
The source reports derive the following candidate specializations of
boundary likelihood experiments. They concern correctly specified iid
sampling with positive nondegenerate face parameters and the face nuisance
parameters fitted. The binomial comparison requires known size at least two;
the NB and unknown-size binomial alternatives to Poisson are tested
separately, rather than as their union. They are not blanket approvals of
the library's tests or confidence intervals.
\begin{center}
\small
\begin{tabular}{p{.37\textwidth}p{.23\textwidth}p{.29\textwidth}}
\toprule
Null versus enlargement&Local geometry&Likelihood-ratio limit\\
\midrule
Exponential versus Lomax; Poisson versus NB or unknown-size binomial;
binomial versus beta-binomial&Gaussian ray&$\tfrac12\chi_0^2+\tfrac12\chi_1^2$\\
Weibull versus Burr&Gaussian ray, efficient information $1-6/\pi^2$&Same mixture\\
Poisson versus beta-binomial with size and both shapes free&Gaussian half-plane&$\tfrac12\chi_1^2+\tfrac12\chi_2^2$\\
Pareto versus Burr with endpoint fitted&Endpoint Poisson process and independent Gaussian tail factor&Point-process supremum\\
\bottomrule
\end{tabular}
\end{center}
The ray mixture is standard once the local geometry and nuisance
regularity are established \cite{self}. The efficient Burr constant
$1-6/\pi^2\approx0.392$ and the beta-binomial half-plane were not located
in the inspected literature. For a fitted Pareto endpoint the $n^{-1}$
endpoint displacement produces a point-process likelihood, so the ordinary
chi-bar-square conclusion cannot be inferred from the parameter being on
a boundary. Individual family assumptions and uniform estimator transfer
remain to be checked.
\todo{Decide which derived closure experiments belong in the main paper. Obtain theorem-level comparisons and verify each family's assumptions before stating inferential calibration as established.}
\source{\code{2026-11-04-infer-closure-report.md}; \code{boundary-points.tex}.}

\section{Software findings and coverage}\label{sec:software}
\subsection{What the inventory counts}
A family-level certification entry aggregates its supported held subsets.
A fully certified entry asserts global reasoning for each such route under
its arithmetic convention; it does not assert that every input produces a
finite fit, that all parameter holds are supported, or that all coefficient
accuracy targets have already been unified. Partial certification means
some routes have certificates and others still have numerical searches.
\begin{center}
\begin{tabular}{lrrr}
\toprule
Inventory, 8 November 2026&Full&Partial&Total\\
\midrule
Showcase snapshot&79&20&99\\
Dossier \code{fit.certified} fields in this worktree&81&18&99\\
\bottomrule
\end{tabular}
\end{center}
Both inventories have zero families with no certified subset. These are
moving counts on the same date, not confidence bounds or two replications.
The second row is a read-only census of 99 canonical dossier front matters;
the first is exactly the supplied showcase statement. The difference is
retained rather than silently rewriting the showcase. The census records
metadata, not an independent audit of every proof and implementation.
\todo{Choose and freeze the release revision for the coverage table, reconcile 79 versus 81 with the showcase, and audit whole-family claims against supported holds and delivered coefficient bounds.}

The architecture keeps the statistical outcome in the core. C++ fitters
report status enums without printing or exceptions; adapters turn them
into public conditions. Structured stop details carry model identities,
finite witnesses and likelihood bounds where available. Jacobians and
coefficient maps must be transported along reciprocal and threshold
reductions. A failed parent fit must remain failed after transformation.

The family's density and its fitted domain are also explicit. Finite proper
discrete point masses can be attained members; continuous zero-scale point
masses are outside the Lebesgue-dominated family. Shape restrictions are
chosen for existence and are flagged at a fitted restriction. A restriction
such as GPD $\xi\ge-1$ does not imply regular inference throughout that
range. Holding a threshold by interface policy is not itself a theorem
that a free threshold never has a maximum.

\subsection{Failure examples and evaluation layers}
The failure gallery reproduces ten textbook likelihood pathologies and
historical implementation failures. They include start values returned as
fits, local maxima selected without global comparison, and false reports
of Poisson or ZIP escape. The near-face sample in
Proposition~\ref{prop:near} shows why a likelihood-gain tolerance proportional
to $n$ is not a statistical nonexistence criterion. The paralogistic and
$t$ examples show why multiple starts and stable Cauchy-like endpoints
do not certify the shape profile.

Density evaluation is an additional dependency. A finite log mass should
remain finite even when the mass itself underflows. The retained
extraDistr 1.10.0.5 reproducer reports \code{-Inf} for
\begin{quote}\small
\verb|extraDistr::dzinb(10000, 2, 0.5, 0.1, log = TRUE)|
\end{quote}
with expected value about $-6923.75$. This is the positional call in the
saved verification script. The installed 1.10.0.5 signature names its third
argument \code{prob}, not \code{mu}. Thus the expected value here uses
base probability $p=0.5$. The showcase names \code{mu} instead; that
transcription is not the retained reproducer. Computing a small probability and then taking its logarithm can produce
this failure even when the true log mass is finite. Such an evaluation
failure is distinct from an optimization failure.

The retained actuar 3.3.7 calculation provides a separate PIG example:
\begin{quote}\small
\verb|dpoisinvgauss(2, mean = 2, shape = 1e20, log = TRUE)|
\end{quote}
returned $0.6931471806$, whereas the reported reference is approximately
$-1.306852819$. A positive log mass is already impossible for a discrete
probability law. The reference is consistent with the Poisson limit at
large shape. The same archive records underflow in logarithmic truncated
count masses and upper-tail CDFs. These are versioned historical findings;
no claim about a later release is implied.
\todo{Choose whether to retain the versioned external-package defect examples. Confirm the archived versions and parameter conventions for publication, and do not imply that a historical defect persists in later releases.}

The validation hierarchy is deliberately heterogeneous. A symbolic-density
comparison can detect a parameter-map error; a high-precision independent
root can check coefficient accuracy; tail truth tables test distribution
functions; regression goldens detect drift. A complete root bracket plus a
uniqueness theorem can prove globality conditional on signs. A compact
cover plus tail inequalities can do the same without uniqueness. No one
layer replaces the others. The mathematical assertions in this paper cite
their reports rather than using the number of passing tests as proof.

Performance evidence is sample-specific. Grouping repeated observations can
reduce the cost of valid bound evaluations without changing the statistical
problem. Conversely a nearly flat likelihood can make certification expensive
even for a short coefficient vector. Budget and numerical stops expose this
cost instead of converting it into an unchecked estimate. A reported
finite witness is useful for diagnosis but must not be silently coerced
into an MLE by a downstream binding.
\source{\code{failure-gallery.tex}; \code{dev/research/upstream/extraDistr.md}, \code{actuar.md} and \code{verify.R}; family certificate reports; read-only dossier census accompanying this draft's report.}

\section{Open problems and discussion}\label{sec:open}
\subsection{Sampford's conjecture: an exact remaining inequality}
Sampford \cite{sampford}, Section 8, left necessity and uniqueness for the
two-parameter truncated-NB likelihood conjectural. A later settlement was
not located in the bounded searches; the source ledger distinguishes later works
actually read from inaccessible leads. That access history does not establish that
the conjecture is still open in all published literature.

Let $P$ be the empirical distribution of positive counts, with mean $m>1$.
At size $r>0$, the conditional probability maximum $p=e^{-t(r)}$ solves
\[
 m=\frac{r(e^t-1)}{1-e^{-rt}}.
\]
Put $h_r(x)=\sum_{i=1}^{x-1}i/(r+i)$,
$b_r(x)=\sum_{i=1}^{x-1}i/(r+i)^2$, $J(t)=1-t/(e^t-1)$, and
$E(u)=1/(e^u-1)-1/u$. For the profile $T$,
\begin{align*}
 F(r):=rT'(r)/k&=mJ(t)-\E_P h_r(X),\\
 t'(r)&=\frac{tE(rt)}{1+E(t)-rE(rt)}<0,\\
 F'(r)&=\E_Pb_r(X)-\mathcal D(r,m),\qquad
 \mathcal D=-mJ'(t)t'(r)>0.
\end{align*}
A sufficient strict stationary-curvature inequality is
\begin{equation}\label{eq:conjecture}
 \E_P h_r(X)=mJ(t)\quad\Longrightarrow\quad
 \E_P b_r(X)<\mathcal D(r,m).
\end{equation}
It would force every stationary point to be a downward crossing, giving
uniqueness and the missing necessity of the admissible ZIP-face condition
$C>0$ in the ZINB excess-zero region. It has not been proved here.

\begin{theorem}[Proved reductions for the ZTNB remainder]\label{thm:reduction}
At a fixed positive mean, truncated-NB laws strictly decrease in convex
order as size grows and dominate the corresponding truncated Poisson.
If the empirical positive-count law is below its fitted truncated Poisson
in convex order, its size profile increases strictly, so no finite-size
maximum exists. This covers every sample on two adjacent positive integers.
Further, maximizing $\E_Pb_r(X)$ with prescribed $\E_PX=m$ and
$\E_Ph_r(X)=H$ reduces exactly to a law supported on $\{1,j,j+1\}$ for
an appropriate $j\ge2$.
\end{theorem}
\begin{proof}[Proof sketch]
For sizes $r_1<r_2$ at fixed mean, the ratio of successive mass ratios
is $(1-e^{-t_2})(r_2+x)/\{(1-e^{-t_1})(r_1+x)\}$, decreasing in $x$.
Equal means and normalization give the two-crossing sign pattern that
implies convex order; the Poisson comparison follows by the limiting
ratio. The score kernel $h_r$ is strictly convex in the integer count.
Its expected value is strictly larger under the finite-size model than
under the fitted Poisson, and no smaller under that Poisson than under
$P$. The score identity then gives $F(r)>0$.

For the sharp reduction, with $\Delta g_x=g_{x+1}-g_x$, direct differences
give
\[
 \frac{\Delta^2 b_r(x)}{\Delta^2 h_r(x)}
   =\frac1{r+x}+\frac1{r+x+1}-\frac1r,
\]
a strictly decreasing sequence. The plane
$b_r(x)\le\alpha_jh_r(x)+\beta_j(x-1)$ interpolating at $1,j,j+1$
therefore dominates at all other positive integers, strictly. Prescribing
$\eta=H/(m-1)$, choose $j$ so that
$h_r(j)/(j-1)\le\eta\le h_r(j+1)/j$. The unique mixture at those three
points with the prescribed moments is nonnegative, by convexity of $h_r$,
and attains the supporting-plane bound. Taking expectations proves the
reduction.
\end{proof}
The remaining scalar inequality is
\[
 \mathcal D(r,m)>\alpha_j(r)mJ(t(r,m))+\beta_j(r)(m-1).
\]
Even its $j=2$ specialization remains unproved over its full feasible
region. The retained calculation found positivity at 1,012 evaluations,
with sizes $10^{-12}$ to $10^{12}$, means from $1.00000001$ to $1000001$,
and extremal indices up to $1,802,147$. This is numerical evidence,
not a cover of a continuous domain. A computational proof would need
validated region bounds and all four limiting regimes. Real-frequency
extremizers also require care: a real-weight equality is not automatically
an empirical rational-frequency counterexample.

The attained size-zero logarithmic-series member must remain in the domain
statement. In the classical positive-size problem there are two endpoint
conditions; in the package's extended nonnegative-size domain, a finite
maximum can instead be the logarithmic member. All-one samples form a
separate point-mass nonidentification case. Confusing these versions would
state the wrong conjecture.
\source{\code{2026-11-07-ztnb-conjecture-report.md}, including supporting-plane proof and source ledger.}

\subsection{Unfinished global covers and arithmetic}
The all-free GEV problem still needs a certificate that closes coupled
shape, location and scale paths. The held-shape Fr\'echet theorem and
held-scale/nonpositive-shape GEV location concavity do not settle that
problem. Partial certificates for beta prime, BNB, GB2, skew-GED and skew-$t$
also leave joint or held-moment profiles outside their proved routes.
These are explicit missing global covers, not excuses to return an
unchecked optimizer iterate. For the skew families, the supplied polynomial
stationarity equations still need multiple-root and critical-component
handling, convergent compact-box bounds, quantitative face-neighborhood
remainders, and globally certified constrained limit-model profilers.
They do not establish all-free uniqueness or a useful worst-case cost bound.

The paralogistic report records a finite-MLE availability limit at
$x=\operatorname{nextafter}(1,0)$ with scale one. A candidate near
$3.6445696825298163\times10^{17}$ has a narrow root enclosure, but the
padded witness lower bound is too weak to close the small-shape tail.
A tighter stable feasible-likelihood bound is needed. The correct present
outcome is numerical failure, despite existence of a finite maximizer.
Scalar BNB fits can likewise have unique mathematical roots too flat for
the current digamma differences to resolve. Stable differences could
improve availability without changing the theorem.

Accuracy certification near a zero coefficient, almost equal maxima, or
a limiting face needs a deliberate metric. A relative error contract
cannot be read literally at zero; an $O(n)$ likelihood allowance cannot
be substituted for coefficient accuracy or for the $O(1)$ local likelihood
scale needed by asymptotic inference. Uniform arithmetic proofs and
outward-rounded implementations are possible further standards, but the
conditional convention here should be stated honestly before either is
attempted.

\subsection{What remains between a fit and an inference procedure}
For GED with location and power jointly fitted at $1<\nu\le2$, the
fitted-point Hessian proof is missing. For nonregular endpoint and cusp
models, localization and optimization error must be controlled on the
local-experiment scale. For moving joins, derivative adapters and
calibration must respect the piecewise model. Uniform results near
transitions, misspecification and bootstrap validity remain separate
questions. The Pareto/Burr point-process supremum still needs certified
critical-value computation with controlled compact-search error, residual
point-process tail error and Monte Carlo quantile uncertainty. Exact
face-selection probabilities require finer analysis than the first local
experiment; a limiting projection atom is not a finite-sample selection law. A complete maximization certificate does not settle them.

The software contract changes which failures are acceptable. Statistical
nonexistence, a proved tie and a well-diagnosed finite-work refusal are
informative results. A plausible unchecked coefficient vector is not.
The strongest reusable tools in these examples are modest ones: an exact
profile, a monotone score, a convex coordinate, a bound for every escape,
and a comparison whose uncertainty is allowed to stop the calculation.
The resulting library is a collection of stated mathematical problems and
verified routes through them, with its remaining gaps visible.
\todo{Before circulation, settle the main-text versus supplement allocation, select a reproducibility release and archive, and review every open mathematical and literature qualification with the authors. No priority claim is proposed.}

\appendix
\section{Source map and reproducibility of the draft}\label{app:sources}
The manuscript is a synthesis of the 8 November showcase and the named
project reports, not a rerun of their mathematical searches. For concrete
samples taken from floating-point fixtures, printed decimals specify their
round-trip binary64 inputs, not exact rational replacements for those inputs;
certificate comparisons inherit the source's arithmetic convention. The
BNB frequency table is a lossless grouping of its frozen 500-count sample.
The two $t$ samples are the retained weighted fixtures. The beta-prime
sample is the exact decimal list in its refined counterexample file.
Analytic ties use the defining equations instead of pretending that
rounded decimal approximations are exact solutions.

The companion draft report records the source-to-section map, every author
decision, verification discrepancies, commands and final outcomes. PDF
build products live under \code{build/paper-draft/}; the source is
\code{dev/research/paper/draft.tex}. Bibliography entries are embedded so
the document compiles without a separate reference database. DOIs are
included where the inspected records supply them; preprints and conference
papers without verified DOIs retain explicit access descriptions.

A final paper should freeze the mathematical sources, versioned code and
certificate outputs together. Historical failure descriptions should stay
attached to their tested software versions. A source ledger should retain
what theorem text was read, rather than promoting a DOI resolution or an
abstract into proof verification. No source in this draft is described as
novel solely because a bounded search did not find a predecessor.

\begin{thebibliography}{99}
\bibitem{anscombe} Anscombe, F. J. (1950). Sampling theory of the negative binomial and logarithmic series distributions. \textit{Biometrika} \textbf{37}, 358--382. \href{https://doi.org/10.1093/biomet/37.3-4.358}{doi:10.1093/biomet/37.3-4.358}.
\bibitem{aragon} Arag\'on, J., Eberly, D. and Eberly, S. (1992). Existence and uniqueness of the maximum likelihood estimator for the two-parameter negative binomial distribution. \textit{Statistics \& Probability Letters} \textbf{15}, 375--379. \href{https://doi.org/10.1016/0167-7152(92)90157-Z}{doi:10.1016/0167-7152(92)90157-Z}. Project proof comparison used the accessible 1991 Stanford preprint, not verified published proof text.
\bibitem{levin} Levin, B. and Reeds, J. (1977). Compound multinomial likelihood functions are unimodal: proof of a conjecture of I. J. Good. \textit{Annals of Statistics} \textbf{5}, 79--87. \href{https://doi.org/10.1214/aos/1176343741}{doi:10.1214/aos/1176343741}.
\bibitem{wang} Wang, Y. (1996). Estimation problems for the two-parameter negative binomial distribution. \textit{Statistics \& Probability Letters}. \href{https://doi.org/10.1016/0167-7152(94)00259-2}{doi:10.1016/0167-7152(94)00259-2}. Abstract and metadata inspected; full objection not acquired.
\bibitem{brons} Br\o ns, H. and Jensen, S. T. (1988). Maximum likelihood estimation in the negative binomial distribution. University of Copenhagen, Institute of Mathematical Statistics, Preprint 1988 No. 8. Complete preprint inspected in the project NB proof report; no DOI verified.
\bibitem{willmot} Willmot, G. E. (1987). The Poisson-inverse Gaussian distribution as an alternative to the negative binomial. \textit{Scandinavian Actuarial Journal}, 113--127. \href{https://doi.org/10.1080/03461238.1987.10413823}{doi:10.1080/03461238.1987.10413823}.
\bibitem{jones} Jones, M. C. (2009). Kumaraswamy's distribution: a beta-type distribution with some tractability advantages. \textit{Statistical Methodology} \textbf{6}, 70--81. \href{https://doi.org/10.1016/j.stamet.2008.04.001}{doi:10.1016/j.stamet.2008.04.001}.
\bibitem{wingo83} Wingo, D. R. (1983). Maximum likelihood methods for fitting the Burr type XII distribution to life test data. \textit{Biometrical Journal} \textbf{25}, 77--84. \href{https://doi.org/10.1002/bimj.19830250109}{doi:10.1002/bimj.19830250109}.
\bibitem{wingo93} Wingo, D. R. (1993). Maximum likelihood estimation of Burr XII distribution parameters under Type II censoring. \textit{Microelectronics Reliability} \textbf{33}, 1251--1257. \href{https://doi.org/10.1016/0026-2714(93)90126-J}{doi:10.1016/0026-2714(93)90126-J}.
\bibitem{grimshaw} Grimshaw, S. D. (1993). Computing maximum likelihood estimates for the generalized Pareto distribution. \textit{Technometrics} \textbf{35}, 185--191. \href{https://doi.org/10.1080/00401706.1993.10485040}{doi:10.1080/00401706.1993.10485040}.
\bibitem{haslett} Haslett, J., Parnell, A. C., Hinde, J. and de Andrade Moral, R. (2022). Modelling excess zeros in count data: a new perspective on modelling approaches. \textit{International Statistical Review}. \href{https://doi.org/10.1111/insr.12479}{doi:10.1111/insr.12479}.
\bibitem{ridout98} Ridout, M., Dem\'etrio, C. G. B. and Hinde, J. (1998). Models for count data with many zeros. International Biometric Conference, Cape Town. Author-hosted conference paper inspected; proceedings pagination and DOI not verified.
\bibitem{ridout01} Ridout, M., Hinde, J. and Dem\'etrio, C. G. B. (2001). A score test for testing a zero-inflated Poisson regression model against zero-inflated negative binomial alternatives. \textit{Biometrics} \textbf{57}, 219--223. \href{https://doi.org/10.1111/j.0006-341X.2001.00219.x}{doi:10.1111/j.0006-341X.2001.00219.x}.
\bibitem{sampford} Sampford, M. R. (1955). The truncated negative binomial distribution. \textit{Biometrika} \textbf{42}, 58--69. \href{https://doi.org/10.1093/biomet/42.1-2.58}{doi:10.1093/biomet/42.1-2.58}.
\bibitem{fernandez} Fern\'andez, C. and Steel, M. F. J. (1999). Multivariate Student-t regression models: pitfalls and inference. \textit{Biometrika} \textbf{86}, 153--167. \href{https://doi.org/10.1093/biomet/86.1.153}{doi:10.1093/biomet/86.1.153}.
\bibitem{reeds} Reeds, J. A. (1985). Asymptotic number of roots of Cauchy location likelihood equations. \textit{Annals of Statistics} \textbf{13}, 775--784. \href{https://doi.org/10.1214/aos/1176349554}{doi:10.1214/aos/1176349554}.
\bibitem{schoups} Schoups, G. and Vrugt, J. A. (2010). A formal likelihood function for parameter and predictive inference of hydrologic models with correlated, heteroscedastic, and non-Gaussian errors. \textit{Water Resources Research} \textbf{46}, W10531. \href{https://doi.org/10.1029/2009WR008933}{doi:10.1029/2009WR008933}.
\bibitem{tulupyev} Tulupyev, A., Suvorova, A., Sousa, J. and Zelterman, D. (2013). Beta prime regression with application to risky behavior frequency screening. \textit{Statistics in Medicine} \textbf{32}, 4044--4056. \href{https://doi.org/10.1002/sim.5820}{doi:10.1002/sim.5820}.
\bibitem{ueki} Ueki, M. (2023; online 2022). Beta-negative binomial nonlinear spatio-temporal random effects modeling of COVID-19 case counts in Japan. \textit{Journal of Applied Statistics} \textbf{50}, 1650--1663. \href{https://doi.org/10.1080/02664763.2022.2064439}{doi:10.1080/02664763.2022.2064439}.
\bibitem{hallin} Hallin, M. and Ley, C. (2012). Skew-symmetric distributions and Fisher information: a tale of two densities. \textit{Bernoulli} \textbf{18}, 747--763. \href{https://doi.org/10.3150/12-BEJ346}{doi:10.3150/12-BEJ346}.
\bibitem{rubio} Rubio, F. J. and Steel, M. F. J. (2014). Inference in two-piece location-scale models with Jeffreys priors. \textit{Bayesian Analysis} \textbf{9}, 1--22. \href{https://doi.org/10.1214/13-BA849}{doi:10.1214/13-BA849}.
\bibitem{smith} Smith, R. L. (1985). Maximum likelihood estimation in a class of nonregular cases. \textit{Biometrika} \textbf{72}, 67--90. \href{https://doi.org/10.1093/biomet/72.1.67}{doi:10.1093/biomet/72.1.67}.
\bibitem{bucher} B\"ucher, A. and Segers, J. (2017). On the maximum likelihood estimator for the generalized extreme-value distribution. \textit{Extremes} \textbf{20}, 839--872. \href{https://doi.org/10.1007/s10687-017-0292-6}{doi:10.1007/s10687-017-0292-6}.
\bibitem{frechet2016} B\"ucher, A. and Segers, J. (2016 manuscript). Maximum likelihood estimation for the Fr\'echet distribution based on block maxima extracted from a time series. \href{https://arxiv.org/abs/1511.07613}{arXiv:1511.07613v2}. Inspected manuscript version; no journal DOI assigned here.
\bibitem{bottazzi} Bottazzi, G. and Secchi, A. (2008 working-paper version, 4 December). Maximum likelihood estimation of the symmetric and asymmetric exponential power distribution. LEM Working Paper 2006/19. Theorems 2.1--2.2 inspected in the project literature ledger; no DOI verified for this version.
\bibitem{cheng} Cheng, R. C. H. and Iles, T. C. (1987). Corrected maximum likelihood in non-regular problems. \textit{Journal of the Royal Statistical Society, Series B} \textbf{49}, 95--101. \href{https://doi.org/10.1111/j.2517-6161.1987.tb01430.x}{doi:10.1111/j.2517-6161.1987.tb01430.x}. Summary and metadata only in the project search.
\bibitem{hallwang} Hall, P. and Wang, J. (2005). Bayesian likelihood methods for estimating the end point of a distribution. \href{https://doi.org/10.1111/j.1467-9868.2005.00523.x}{doi:10.1111/j.1467-9868.2005.00523.x}. Abstract and metadata only in the project search.
\bibitem{self} Self, S. G. and Liang, K.-Y. (1987). Asymptotic properties of maximum likelihood estimators and likelihood ratio tests under nonstandard conditions. \textit{Journal of the American Statistical Association} \textbf{82}, 605--610. \href{https://doi.org/10.1080/01621459.1987.10478472}{doi:10.1080/01621459.1987.10478472}.
\end{thebibliography}
\end{document}
